\section*{الفصل \ref{ch:b2:diatomic-mos} --- المدارات الجزيئية للجزيئات الثنائية الذرة}

\begin{solution}{exo:b2:diatomic-mos:1}
$E_+ = (-13.6 - 9.0)/1.60 = \qty{-14.1}{eV}$؛ $E_- = (-13.6 + 9.0)/0.40 = \qty{-11.5}{eV}$.
إلكترونان في $\psi_+$: $-28.25$ مقابل $2 \times (-13.6) = \qty{-27.2}{eV}$، أي ارتباط قدره
\qty{1.05}{eV} في هذا النموذج التقريبي.
\end{solution}

\begin{solution}{exo:b2:diatomic-mos:2}
\ce{He2}: $(2 - 2)/2 = 0$؛ \ce{He2+}: $(2 - 1)/2 = 0.5$؛ \ce{Li2}: 1؛ \ce{B2}: ستة إلكترونات
تكافؤ، $\sigma_{2s}^2\sigma_{2s}^{*2}\pi_{2p}^2$، \omterm{def:b2:diatomic-mos:bond-order}{رتبة الرابطة} 1.
\end{solution}

\begin{solution}{exo:b2:diatomic-mos:3}
\ce{B2} (إلكترونان منفردان في المدارين $\pi_{2p}$، مع مزج s--p) 
و\ce{O2} (إلكترونان في $\pi^*$). أما \ce{C2} و\ce{N2} و\ce{F2} \omterm{def:b2:diatomic-mos:magnetism}{فدايامغناطيسية}.
\end{solution}

\begin{solution}{exo:b2:diatomic-mos:4}
غير معدوم: (1s، $2\mathrm p_z$)، و($2\mathrm p_x$، $2\mathrm p_x$)، و(2s، $2\mathrm p_z$).
ومعدوم بسبب التناظر: (1s، $2\mathrm p_x$) و($2\mathrm p_x$، $2\mathrm p_y$).
\end{solution}

\begin{solution}{exo:b2:diatomic-mos:5}
ينزع التأين إلكترونًا رابطًا: \omterm{def:b2:diatomic-mos:bond-order}{رتبة الرابطة} 2.5 بدل 3، فتكون الرابطة أطول
وأضعف، كما يبيّن $D_0(\ce{N2+}) = \qty{8.71}{eV} < D_0(\ce{N2}) = \qty{9.76}{eV}$.
\end{solution}

\begin{solution}{exo:b2:diatomic-mos:6}
تقع مدارات الأكسجين أدنى (فالأكسجين أكثر كهرسلبية): وحسب
\cref{thm:b2:diatomic-mos:heteronuclear}، تكون \omterm{def:b2:diatomic-mos:bonding}{المدارات الرابطة} مرجَّحة على
الذرة الأدنى، الأكسجين، والمدارات الأعلى على الكربون. وأعلى مدار
مشغول، وهو مدار $\sigma$ ذو طابع كربوني غالب يتجه بعيدًا عن الأكسجين، هو
الزوج الحر على الكربون.
\end{solution}

\begin{solution}{exo:b2:diatomic-mos:7}
لا يملك تناظر المدار 1s للهيدروجين إلا المدار $2\mathrm p_z$ للفلور (على طول المحور):
فيكوّنان مدارًا $\sigma$ مرجَّحًا على F ومدارًا $\sigma^*$ مرجَّحًا على H. أما
$2\mathrm p_x$ و$2\mathrm p_y$ للفلور (تداخل معدوم) ومداره 2s المنخفض (البعيد جدًا طاقيًا)
فتبقى غير رابطة: ثلاثة أزواج حرة. وتملأ إلكترونات التكافؤ الثمانية 2s و$\sigma$
والمدارين p غير الرابطين.
\end{solution}

\begin{solution}{exo:b2:diatomic-mos:8}
$\bar\alpha = \qty{-12}{eV}$، $\Delta = \qty{4}{eV}$، $\sqrt{4 + 9} = \qty{3.61}{eV}$:
$E = \qty{-15.6}{eV}$ و\qty{-8.4}{eV}. ومن أجل المستوى الأدنى $(\alpha_B - E)c_B = -\beta
c_A$: $1.61c_B = 3.0c_A$، $c_A/c_B = 0.535$، $c_B^2 = 1/(1 + 0.286) = 0.78$.
\end{solution}

\begin{solution}{exo:b2:diatomic-mos:9}
ثمانية إلكترونات تكافؤ: $\sigma_{2s}^2\sigma_{2s}^{*2}\pi_{2p}^4$، والمدار $\sigma_{2p}$
(المدفوع فوق $\pi_{2p}$ بالمزج) فارغ. \omterm{def:b2:diatomic-mos:bond-order}{رتبة الرابطة} $(6 - 2)/2 = 2$، من
مداري $\pi$ الممتلئين.
\end{solution}

\begin{solution}{exo:b2:diatomic-mos:10}
NO، أحد عشر إلكترون تكافؤ، أحدها في $\pi^*$: \omterm{def:b2:diatomic-mos:bond-order}{رتبة الرابطة} 2.5. \ce{NO+}، عشرة: \omterm{def:b2:diatomic-mos:bond-order}{رتبة الرابطة}
3. $D_0(\ce{NO+}) = 6.51 + 13.62 - 9.26 = \qty{10.87}{eV}$: أقوى بكثير من NO، لأن
الإلكترون المنزوع كان مضادًا للربط.
\end{solution}

\begin{solution}{exo:b2:diatomic-mos:11}
$E_+ + E_- = \dfrac{(\alpha + \beta)(1 - S) + (\alpha - \beta)(1 + S)}{1 - S^2} = \dfrac{2(\alpha -
\beta S)}{1 - S^2}$. ويتجاوز $2\alpha$ حين $\alpha - \beta S > \alpha - \alpha S^2$، أي
$\beta < \alpha S$ (بالقسمة على $-S < 0$). ومع أربعة إلكترونات، اثنان في كل مدار، تكون
الطاقة الكلية فوق طاقة المدارين المنفصلين: فالمداران الممتلئان يتنافران.
\end{solution}

\begin{solution}{exo:b2:diatomic-mos:12}
رتب الرابطة 2.5 و2 و1.5 و1: وبهذا الترتيب تطول الروابط وتضعف. ويحوي بيروكسيد
الهيدروجين الرابطة الأحادية O--O للبيروكسيد (الوحدة \ce{O2^2-})؛ ويحوي \ce{KO2}
أيون الأكسيد الفائق \ce{O2-}، \omterm{def:b2:diatomic-mos:magnetism}{البارامغناطيسي}.
\end{solution}

\begin{solution}{pb:b2:diatomic-mos:1}
\textbf{1.} اثنا عشر.
\textbf{2.} $\sigma_{2s}^2\,\sigma_{2s}^{*2}\,\sigma_{2p}^2\,\pi_{2p}^4\,\pi_{2p}^{*2}$، وإلكترونا $\pi^*$
في مدارين مختلفين.
\textbf{3.} $(8 - 4)/2 = 2$.
\textbf{4.} اثنان: \ce{O2} \omterm{def:b2:diatomic-mos:magnetism}{بارامغناطيسي}.
\textbf{5.} إنها تقرن كل الإلكترونات، بينما يقع اثنان منها منفردين في مداري
$\pi^*$ المتساويي الطاقة.
\textbf{6.} $D_0 = 2 \times 246.79 = \qty{493.6}{kJ/mol}$، أي \qty{5.12}{eV}.
\textbf{7.} \ce{O2+}: $\pi^{*1}$، \omterm{def:b2:diatomic-mos:bond-order}{رتبة الرابطة} 2.5.
\textbf{8.} \ce{O2-}: $\pi^{*3}$، \omterm{def:b2:diatomic-mos:bond-order}{رتبة الرابطة} 1.5.
\textbf{9.} \ce{O2^2-}: $\pi^{*4}$، \omterm{def:b2:diatomic-mos:bond-order}{رتبة الرابطة} 1.
\textbf{10.} \ce{O2+} واحد، و\ce{O2-} واحد، و\ce{O2^2-} لا شيء.
\textbf{11.} \babelsublr{\ce{O2+} $<$ \ce{O2} $<$ \ce{O2-} $<$ \ce{O2^2-}}.
\textbf{12.} أيون البيروكسيد.
\textbf{13.} \ce{O2+} و\ce{O2} و\ce{O2-}.
\textbf{14.} من مدار $\pi^*$.
\textbf{15.} المدار $\pi^*$، المضاد للربط، يقع فوق المستوى 2p للذرة: فيكون
إلكترونه أقل ارتباطًا.
\textbf{16.} مساران من \ce{O2} إلى $\ce{O} + \ce{O+} + \mathrm e^-$: التفكيك ثم
تأيين ذرة، $D_0(\ce{O2}) + I(\ce{O})$؛ أو تأيين الجزيء ثم تفكيك الأيون،
$I(\ce{O2}) + D_0(\ce{O2+})$. وهما متساويان: $D_0(\ce{O2+}) = 5.12 + 13.62 - 12.07 = \qty{6.66}{eV}$.
\textbf{17.} أكبر من $D_0(\ce{O2}) = \qty{5.12}{eV}$: فنزع إلكترون مضاد للربط
يقوي الرابطة (الرتبة 2.5).
\textbf{18.} $D_0(\ce{N2}) = 2 \times 470.82/96.485 = \qty{9.76}{eV}$؛ $D_0(\ce{N2+}) = 9.76 +
14.53 - 15.58 = \qty{8.71}{eV}$.
\textbf{19.} أعلى مدار مشغول في \ce{N2} رابط، يقع تحت المستوى 2p
للذرة.
\textbf{20.} حين يأتي الإلكترون المنزوع من \omterm{def:b2:diatomic-mos:bonding}{مدار مضاد للربط}.
\textbf{21.} قاعدة هوند (أكبر عدد من السبينات المتوازية في المدارات المتساوية الطاقة).
\textbf{22.} إلكترونا $\pi^*$ كلاهما في المدار نفسه، بسبينين مقترنين.
\textbf{23.} يتنافر إلكترونان في المدار نفسه أكثر، ويفقدان
استقرار التبادل الخاص بالسبينات المتوازية.
\textbf{24.} تبقى \omterm{def:b2:diatomic-mos:bond-order}{رتبة الرابطة} 2؛ وكل الإلكترونات مقترنة: \omterm{def:b2:diatomic-mos:magnetism}{دايامغناطيسي}.
\textbf{25.} لأيون \ce{O2+} رتبة رابطة $\boldsymbol{2.5}$ وطاقة تفكك
$\boldsymbol{\approx \qty{6.66}{eV}}$، مقابل 2 و\qty{5.12}{eV} للجزيء \ce{O2}.
\end{solution}
